\section{Address Mapping Reconstruction} \label{sec:targeted}
    After bypassing the cache hierarchy, a fundamental challenge remains, namely the existence of a 
    fixed mapping between physical addresses and the \gls{dram} layout, as noted in~\cite{toolRPI3}. 
    To achieve "\textit{\hyperref[ch:targeted]{C2. Address Mapping Reconstruction}}", it is necessary to access a specific 
    \gls{dram} row within the same bank in order to gain bit flips in an adjacent victim row through produced row conflicts. 
    However, the exact informations of this physical-to-\gls{dram} address mapping are typically not publicly available.
    While AMD provided documentation of the DRAM bank addressing functions for older processor models - specifically in 
    the BIOS and Kernel Developer's Guide (BKDG) up to microarchitecture 16h~\cite{webAMDDeveloperCentral} - they ceased 
    to publish this information with the beginning of microarchitecture 17h (released in 2017~\cite{webZenCPU})~\cite{paperREBankAddressing}.
    Unlike AMD's earlier documentation efforts, Intel has never disclosed the DRAM addressing functions implemented in its 
    processors~\cite{paperREBankAddressing},which is also the case for all our ARM test devices used in this work.
    Consequently, the only viable approach is to 
    reverse engineer the \gls{dram} address mapping. This section briefly discusses the available methods, 
    explains why certain approaches are adopted or discarded, and identifies the most 
    effective technique for ARM-based systems such as the Raspberry Pi. In addition, it outlines the 
    practical implementation process, including the development steps and the challenges 
    encountered during implementation.

    \subsection{Methods} \label{subsec:methods}
        If the internal organization of \gls{dram} and the behavior of the row buffer are known, it becomes 
        feasible to identify two distinct rows located within the same bank. Based on this observation, 
        several approaches have been proposed to reverse engineer the mapping between physical addresses 
        and \gls{dram} structures. These approaches mainly differ in how address sets are constructed and how 
        the underlying addressing functions are inferred. This work focuses on three representative methods, which
        are listed below, along with a brief description of each.
        \\ \\
        Pessl et al.~\cite{paperDRAMA} in 2016, exploits a row-buffer timing side channel in which 
        row conflicts manifest as increased memory access latency. By observing row hits and row conflicts, physical 
        addresses can be grouped into sets that reside within the same bank and thus share identical channel, DIMM, 
        rank, and bank properties. The underlying \gls{dram} address mapping is inferred by exhaustively evaluating 
        linear addressing functions that distinguish these sets, which are commonly expressed as XOR combinations of 
        physical address bits. This general methodology underlies several subsequent attacks and tools, 
        including~\cite{paperHalfDouble, toolHalfDouble, paperDrammer, toolDrammer, thesisCacheAttacksandRH}.
        \\
        The work introduced by Xiao et al.~\cite{paperOneBitFlipsOneCouldFlops} in 2016, proposes a 
        graph-based approach to infer physical-to-\gls{dram} address mappings. Similar to DRAMA, it relies on 
        timing differences caused by row conflicts, but its primarily aimed at enableing Rowhammer attacks 
        across virtual machines.
        \\
        Plin et al.~\cite{paperKnockKnock} in 2025, reformulates the identification of \gls{dram} parity 
        masks as a linear algebra problem as Helm et al.~\cite{paperReliableReverseEngineering} in 2020. Instead of relying on exhaustive search, this method focuses on 
        solving the address mapping directly through algebraic reasoning.
        \\
        The following sections assess the suitability of these methods by applying them to the test systems and confirming 
        or refuting them based on experimental results.

    \input{pa/reverse-engineering}

    \input{pa/verification}
